Los proyectos residenciales de empresas desarrolladoras en la Ciudad de Guatemala generan un volumen de leads que excede la capacidad de atención de la fuerza de ventas, sin un criterio estructurado que oriente sobre a quién contactar primero. Bajo una selección aleatoria, se necesita contactar, en promedio, a 31 leads para obtener una conversión, cifra que refleja el esfuerzo comercial asociado a una atención sin priorización.

Este trabajo desarrolló un sistema de priorización comercial compuesto por tres elementos: un modelo que estima la probabilidad de conversión de cada lead, una segmentación que perfila a los clientes potenciales y un tablero en Power BI que presenta ambos resultados al equipo comercial. El proceso de ventas se documentó previamente mediante notación BPMN para identificar el punto del flujo en el que debía incorporarse la priorización.

El resultado central es una reducción sustancial del esfuerzo comercial por conversión. Al priorizar la cartera mediante el modelo, el 5\% con mayor probabilidad estimada concentra 223 de las 739 conversiones históricas del universo analizado, lo que equivale a 5 leads contactados por conversión obtenida, frente a los 31 de la selección aleatoria. En un punto de operación más amplio, el modelo recupera el 70 \% de las conversiones reales al priorizar menos de un tercio de la cartera. La segmentación identificó seis perfiles de clientes, dos de los cuales presentan una tasa de conversión significativamente superior al promedio y comparten una misma motivación de compra.

El sistema reordena la cartera sin garantizar la identificación de compradores individuales; por lo tanto, se recomienda usarlo como criterio de orden de atención, reforzar el registro de las variables clave y validar su impacto mediante una prueba controlada antes de su adopción generalizada.
